\section{Introduction}

Myelinating oligodendrocytes (OLs) in the central nervous system do more than accelerate action potential propagation: they also provide trophic and metabolic support to the axons they ensheathe \citep{Nave2010,Simons2016,Philips2017}. A seminal observation that myelin supplies lactate to energy-demanding axons through monocarboxylate transporter 1 (MCT1) established a bidirectional, activity-dependent dialogue between axons and their glial partners \citep{Funfschilling2012,Lee2012}. Subsequent work showed that oligodendroglial N-methyl-D-aspartate (NMDA) receptors regulate GLUT1-mediated glucose import into myelin and thereby tune energy delivery to active axons \citep{Saab2016}. Together these studies redefined CNS myelin as a metabolic compartment whose integrity depends on an ongoing conversation with the underlying axon \citep{Brown2012,Trevisiol2017,Chamberlain2021}.

Despite these advances, the acute signals by which mature OLs sense ongoing axonal firing remain incompletely understood. Candidate messengers include vesicular glutamate acting on oligodendroglial NMDA and AMPA receptors \citep{Micu2006,Saab2016}, ATP acting on purinergic P2X/P2Y receptors \citep{Kettenmann2013}, and, at earlier developmental stages, activity-dependent modulation of proliferation and myelination \citep{Baraban2016,Suminaite2019}. However, each action potential also releases potassium (K$^+$) into the periaxonal space, and this transient extracellular K$^+$ rise has long been suspected to be a fast, energy-efficient carrier of the axonal firing state to surrounding glia \citep{Kofuji2000,Bay2012,Olsen2015}. Inwardly rectifying K$^+$ channels of the Kir4.1 (KCNJ10) family are enriched on oligodendrocytes and astrocytes \citep{Neusch2001,Djukic2007,Larson2018} and are essential for K$^+$ siphoning in white matter \citep{Larson2018,Tong2014}. Conditional loss of oligodendroglial Kir4.1 produces late-onset axonopathy and subtle myelin defects \citep{Schirmer2018,Stumpf2019,Battefeld2019}, but whether Kir4.1 actively couples axonal activity to oligodendroglial signalling and metabolic output in real time has not been directly tested.

Here we combine two-photon Ca$^{2+}$ imaging in OLs of acutely isolated optic nerves with simultaneous compound action potential (CAP) recordings, pharmacology, OL-specific Kir4.1 knockout, and genetically encoded metabolite sensors to address this question. We find that mature OLs respond to high-frequency axonal firing with robust Ca$^{2+}$ surges whose amplitude scales with stimulation frequency and that depend on action potentials, extracellular Ca$^{2+}$, and oligodendroglial Kir channels, with reverse-mode Na$^+$/Ca$^{2+}$ exchanger (NCX) activity as the dominant Ca$^{2+}$ entry pathway \citep{Stys1992,Tekkok2005,Micu2018}. Oligodendrocyte-specific deletion of Kir4.1 slows white-matter K$^+$ clearance, produces late-onset axonal pathology, and decreases the abundance of MCT1 and GLUT1 in central myelin, leading to reduced axonal lactate availability and blunted activity-induced glucose consumption. These findings establish extracellular K$^+$ and oligodendroglial Kir4.1 as a fast axon-to-OL signalling axis that gates the metabolic support required to sustain long-term axonal integrity.

\section{Related Work}

\subsection{Oligodendrocyte-to-axon metabolic coupling}

Genetic strategies that eliminate myelin-specific monocarboxylate transporter MCT1 or shift myelin toward glycolysis revealed that OLs export lactate to fuel axons whose electrical activity outpaces local mitochondrial capacity \citep{Funfschilling2012,Lee2012}. Rinholm and colleagues showed that glucose and lactate supply alone are sufficient to modulate OL differentiation and myelination in vitro \citep{Rinholm2011}. Saab et al.\ subsequently demonstrated that oligodendroglial NMDA receptors trigger surface expression of GLUT1, linking glutamatergic signalling to metabolite transport \citep{Saab2016}. Trevisiol and colleagues used ATeam-based FRET imaging to show in real time that axonal ATP in white matter depends on glycolytic activity in surrounding OLs \citep{Trevisiol2017}. These studies establish a framework in which axonal activity drives OL metabolism, but they leave open the identity of the earliest, obligatory activity signal that reaches mature OLs.

\subsection{Potassium dynamics and Kir4.1 in glia}

Kir4.1 is the principal inwardly rectifying K$^+$ channel in macroglia and sets a hyperpolarized membrane potential that both follows and buffers extracellular K$^+$ transients \citep{Kofuji2000,Olsen2015,Hawkins2017}. Genetic removal of Kir4.1 from astrocytes or oligodendrocyte lineage cells impairs K$^+$ and glutamate uptake, causes early lethality in global knockouts, and produces progressive white-matter defects when restricted to oligodendrocytes \citep{Djukic2007,Neusch2001,Larson2018,Schirmer2018}. Pharmacological blockade of Kir channels with barium slows K$^+$ clearance in optic nerve and alters axonal conduction recovery after high-frequency firing \citep{Bay2012,Larson2018}. Loss-of-function of Kir4.1 is also implicated in Huntington's disease and multiple sclerosis pathology, placing this channel at the intersection of homeostasis and neurological disease \citep{Tong2014,Schirmer2018,Stumpf2019}.

\subsection{Ionic mechanisms of OL Ca$^{2+}$ signalling}

Myelinic and somatic Ca$^{2+}$ responses in OLs have been attributed to glutamate receptors \citep{Micu2006,Saab2016}, purinergic receptors \citep{Kettenmann2013}, and voltage-gated Ca$^{2+}$ channels. Reverse-mode operation of the Na$^+$/Ca$^{2+}$ exchanger (NCX) has been documented in CNS white matter where it contributes to anoxic Ca$^{2+}$ overload and axonal injury \citep{Stys1992}. More recent axo-myelinic signalling models propose local bidirectional communication at the periaxonal space \citep{Micu2018}. However, the relative contribution of K$^+$/Kir4.1 depolarization versus neurotransmitter-driven mechanisms to acute OL Ca$^{2+}$ rises during physiological firing has not been quantified in intact white matter, motivating the present study.
